% ============================================================
% Abstract (Greek)
%
% Requires Greek language support (texlive-collection-langgreek and
% texlive-babel-english on Fedora). The active document body is
% English; we switch into Greek for this file only and back at the
% end. Latin tokens (acronyms, code identifiers) are wrapped in
% \textlatin{...} so the LGR encoding does not transliterate them.
% ============================================================
\selectlanguage{greek}
\chapter*{Περίληψη}
\addcontentsline{toc}{chapter}{Περίληψη}

\newcommand{\la}[1]{\textlatin{#1}}

Η σύγχρονη πρακτική της διασφάλισης ποιότητας \la{API} (\la{API quality assurance}) είναι κατακερματισμένη σε χαλαρά συνδεδεμένα μεταξύ τους αντικείμενα: τα εργαλεία στατικής ανάλυσης δομής (\la{structural validation linters}) ελέγχουν εάν ένα έγγραφο \la{OpenAPI} αναλύεται συντακτικά, τα εργαλεία δοκιμών συμβολαίου (\la{contract-testing tools}) επαληθεύουν ότι ένας διακομιστής ανταποκρίνεται στις προδιαγραφές του σε μερικά επιλεγμένα τελικά σημεία (\la{endpoints}), και τα συστήματα δοκιμών φόρτου (\la{load-testing harnesses}) μετρούν την απόδοση υπό χειροκίνητα σχεδιασμένα σενάρια. Κάθε εργαλείο είναι ανεξάρτητα χρήσιμο· κανένα όμως δεν αντιμετωπίζει το βαθύτερο πρόβλημα, το οποίο είναι ότι η ίδια η προδιαγραφή (\la{specification}) αποτελεί το μόνο κοινό τεχνούργημα (\la{artefact}) μεταξύ της ανάπτυξης, των δοκιμών, της διάθεσης και της λειτουργίας, και ότι η ποιότητά της καθορίζει την ποιότητα κάθε επόμενης αυτοματοποίησης που εξαρτάται από αυτήν. Η παρούσα διατριβή προτείνει μια διαφορετική αφετηρία.

Οι \textbf{Δοκιμές Βάσει Προδιαγραφών} (\la{Specification-Driven Testing}, \la{SDT}) είναι ένα εννοιολογικό πλαίσιο για την αυτοματοποιημένη διασφάλιση ποιότητας \la{API}, στο οποίο η προδιαγραφή αντιμετωπίζεται ως η κανονική πηγή αλήθειας (\la{canonical source of truth}) τόσο για την \emph{επικύρωση} (περιγράφει το συμβόλαιο ένα καλοσχηματισμένο, πλήρες και ασφαλές \la{API};) όσο και για τις \emph{δοκιμές} (ικανοποιεί η υλοποίηση το συμβόλαιο υπό ρεαλιστικές λειτουργικές συνθήκες και συνθήκες φόρτου;). Το πλαίσιο βασίζεται σε τρία αξιώματα---\emph{Πληρότητα}, \emph{Ντετερμινισμός} και \emph{Παρατηρησιμότητα}---και αναλύεται σε εννέα αρχές (\la{P001} έως \la{P009}) που καλύπτουν τη δομική συμμόρφωση, την ποιότητα της τεκμηρίωσης, τη διαχείριση σφαλμάτων, την ακρίβεια του σχήματος (\la{schema precision}), την ασφάλεια, τη λειτουργική ορθότητα, την απόδοση, τη διαχείριση εκδόσεων (\la{versioning}) και την ετοιμότητα για δοκιμές. Η συνεισφορά σε αυτό το επίπεδο είναι θεωρητική: ένα πλαίσιο, όχι μια πλήρης μεθοδολογία· τα στοιχεία της μεθοδολογίας που απουσιάζουν (τυπική διαδικασία παραγωγής δοκιμών, σημασιολογία εκτέλεσης και αποτυχίας, διακυβέρνηση κύκλου ζωής) ορίζονται ως μελλοντική εργασία.

\textbf{Το \la{DriveBy}} αποτελεί την υλοποίηση αναφοράς (\la{reference implementation}) του πλαισίου: είναι ένας ελεγκτής επικύρωσης γραμμής εντολών (\la{command-line validator}) σε \la{Go}, ο οποίος φορτώνει μια προδιαγραφή \la{OpenAPI 3.0/3.1}, εκτελεί τους επτά επί του παρόντος υλοποιημένους ελεγκτές αρχών (\la{P001--P005, P008, P009}· οι \la{P006} και \la{P007} περιορίζονται στο επίπεδο δοκιμών χρόνου εκτέλεσης) και παράγει μια ντετερμινιστική, αναγνώσιμη από μηχανές αναφορά \la{JSON} με ευρήματα ταξινομημένα ανά σοβαρότητα. Το \la{CLI} είναι δομημένο γύρω από έναν προσαρμογέα \la{APISpec} (\la{APISpec adapter}) που αφαιρεί τις λεπτομέρειες που αφορούν συγκεκριμένες εκδόσεις, επιτρέποντας σε κάθε ελεγκτή να λειτουργεί σε μια ενιαία σημασιολογική επιφάνεια, ανεξάρτητα από το αν η υποκείμενη προδιαγραφή είναι \la{OpenAPI 3.0.x}, \la{3.1.0} ή ένα μετατραμμένο έγγραφο \la{Swagger 2.0}.

\textbf{Το \la{XSDLC}} (\la{Crossplane Software Development Lifecycle composition}) αποτελεί την ενσωμάτωση \la{GitOps}: ένας ενιαίος προσαρμοσμένος πόρος του \la{Kubernetes} (\la{Kubernetes custom resource}) που παρέχει δηλωτικά (\la{declaratively provisions}) μια πλήρη ροή παράδοσης \la{GitOps} δύο αποθετηρίων (\la{repositories})---αποθετήριο εφαρμογής, αποκλειστικό αποθετήριο \la{GitOps}, εφαρμογές \la{ArgoCD} με επικαλύψεις ενυδάτωσης πηγής (\la{source-hydration overlays}), τον \la{GitOps Promoter}, πύλες ποιότητας (\la{quality gates}) πολλαπλών σταδίων που εκτελούν το \la{DriveBy} στα όρια δοκιμών (\la{staging}) και παραγωγής (\la{production}), προστασία κλάδων (\la{branch protection}), καθώς και τις απαραίτητες υποδομές \la{ingress} και πιστοποιητικών για τη λειτουργία του συστήματος. Το \la{XSDLC} είναι το εμπειρικό τεχνούργημα που αποδεικνύει ότι το \la{SDT} μπορεί να ενσωματωθεί ως εγγενές χαρακτηριστικό μιας Εσωτερικής Πλατφόρμας Προγραμματιστών (\la{Internal Developer Platform}) αντί για ένα πρόσθετο εργαλείο στατικής ανάλυσης (\la{bolt-on linter}).

Το πλαίσιο, το \la{CLI} και η σύνθεση (\la{composition}) αξιολογούνται εμπειρικά μέσω τριών σκελών. Μια \emph{ελεγχόμενη αξιολόγηση} σε ένα χειροποίητο \la{API} αναφοράς (\la{non-critical-api}) επαληθεύει την ανίχνευση μεμονωμένων ελαττωμάτων (15 μεταλλάξεις σε επίπεδο ελέγχου, 80\% ποσοστό ανίχνευσης· ανιχνεύθηκαν και οι 3 από τους 3 τυχαίους συνδυασμούς πολλαπλών αρχών). Μια \emph{αξιολόγηση μεγάλης κλίμακας} σε 50 προδιαγραφές \la{OpenAPI 3.x} που συλλέχθηκαν από το μητρώο \la{APIs.guru} δείχνει ότι το 84\% των δημόσιων \la{API} βαθμολογείται με $1/6$ ή χειρότερα στην αυστηρή λειτουργία (\la{strict mode}) του \la{DriveBy}---κανένα \la{API} στο δείγμα δεν περνά περισσότερες από 2 από τις 6 αξιολογημένες αρχές. Μια \emph{επιχειρησιακή απόδειξη ορθότητας εννοίας} (\la{operational proof-of-concept}) αναπτύσσει το \la{XSDLC} στο εσωτερικό \la{cluster Kubernetes} της \la{Novelcore} και επεξεργάζεται πέντε \la{API} μέσω της πλήρους ροής παράδοσης. Ένα μικρό \emph{πείραμα ανατροφοδότησης πρακτόρων} (\la{agent-feedback experiment}) στην κανονική προδιαγραφή \la{Swagger Petstore} δείχνει ότι ένας πράκτορας κώδικα τεχνητής νοημοσύνης (\la{AI coding agent}), λαμβάνοντας ως ανατροφοδότηση μόνο την αναφορά του \la{DriveBy}, μπορεί να αναθεωρήσει αυτόνομα μια προδιαγραφή με βαθμολογία 1/6 σε 4/6 σε μία μόνο επανάληψη---εμπειρική απόδειξη για την τοποθέτηση του πλαισίου ως υποδομή ανατροφοδότησης για τη βελτίωση προδιαγραφών καθοδηγούμενη από πράκτορες (\la{agent-driven specification improvement}).

Η διατριβή κλείνει με μια ρητή αναφορά στην πατρότητα (\la{authorship}) και τη μεθοδολογία ανάπτυξης, διαχωρίζοντας την πνευματική ιδιοκτησία του συγγραφέα όσον αφορά το πλαίσιο, την αρχιτεκτονική και τον σχεδιασμό της αξιολόγησης, από τη φάση παραγωγικοποίησης (\la{productisation}) με τη βοήθεια τεχνητής νοημοσύνης, παρουσιάζοντας ποσοτικά στοιχεία που προκύπτουν από το ιστορικό ελέγχου εκδόσεων (\la{version-control history}).

\noindent\textbf{Λέξεις-κλειδιά:} διασφάλιση ποιότητας \la{API}, δοκιμές βάσει προδιαγραφών, \la{OpenAPI}, \la{Crossplane}, \la{GitOps}, \la{Kubernetes}, πύλες ποιότητας, ανάπτυξη καθοδηγούμενη από πράκτορες.
\selectlanguage{english}

\clearpage
